%------------------------------------------------------------------------------
% \PART 3
%------------------------------------------------------------------------------
\chapter[Разработка нейронной сети для распознавания рукописных цифр]
{РАЗРАБОТКА НЕЙРОННОЙ СЕТИ ДЛЯ РАСПОЗНАВАНИЯ РУКОПИСНЫХ ЦИФР}

%------------------------------------------------------------------------------
% \PART 3.3
%------------------------------------------------------------------------------
\section{Принцип работы LST преобразования}\par
\hspace*{12.5 mm}Двумерное разделимое преобразование используется в обработке 
изображений для уменьшения вычислительной сложности. Основная идея заключается 
в том, что вместо хранения полного двумерного ядра свёртки размером 
n\(\times\)n используется его представление в виде произведения двух одномерных 
векторов. Это позволяет сократить число параметров с \(n^2\) до \(2n\), что 
снижает затраты на вычисления и память.

\begin{equation}
    W = V \times h^T,
\end{equation}

\noindentгде \(W\in \mathbb{R}^{n\times n}\);
\(v, h^T\in \mathbb{R}^{n\times 1}\).

LST основан на идее совместного использования весов одного полносвязанного 
слоя для обработки всех строк изображения. После этого второй общий 
полносвязный слой используется для обработки всех столбцов представления 
изображения, полученных из первого слоя. Использование слоев LST в архитектуре 
нейронной сети значительно сокращает количество параметров модели ($ N = 2 \cdot (d_{in}+1) \cdot d_{out}$) по сравнению 
с моделями, которые используют стекированные полносвязные слои.

Принцип работы LST2D представлен на рисунке~\ref{fig:lst2d-pic}\cite{LST2D}. 

\insertfigure{lst2d-pic}{pic/LST2D.jpg}{Принцип работы LST2D}{16cm}

С математической точки зрения LST2D принимает двумерное изображение X размером 
\(d_{in}\times d_{in}\) в качестве входных данных и создает двумерное выходное 
изображение Y размером $d_{out}$\(\times\)$d_{out}$

\begin{equation}
  Y = LST_{d_{in}\times d_{out}}(X) = \tan(W_2\tan(W_1X^T)).
\end{equation}

\noindentгде W1, W2 — матрицы весов слоев FC1 и FC2 соответственно.

Алгоритм работы LST слоя, показывающий последовательность действий по преобразованию
входного изображения размером \(28\times28\) в выходное, такого же размера 
представлен на рисунке~\ref{fig:alg-lst}. 

\insertfigure{alg-lst}{pic/LST_alg.png}{Алгоритм работы LST2D слоя}{16cm}

Простейшая архитектура сети на основе LST2D показана на 
рисунке~\ref{fig:lst-pic}\cite{LST2D}. 

\insertfigure{lst-pic}{pic/LST.png}{Принцип работы LST-1}{11.0cm}

Данную архитектуру можно масштабировать(рисунок~\ref{fig:lst2-pic}\cite{LST2D}).

\insertfigure{lst2-pic}{pic/LST2.png}{Принцип работы LST-2}{12.5cm}

Пример обработки изображения нейронной сетью архитектуры LST-1 представлен на 
рисунке~\ref{fig:img-proc}\cite{LST2D}.

\insertfigure{img-proc}{pic/img_proc.png}{Архитектура CNN-Transformer}{13cm}

%------------------------------------------------------------------------------
% \PART 3.4
%------------------------------------------------------------------------------
\section{Обучение нейронной сети}
\hspace*{12.5 mm}Обучение нейронной сети LST-1 на основе 60000 тренировочных 
изображений из базы данных MNIST.\@ Для обучения используется фреймворк 
PyTorch.\@

В качестве исходных данных используется набор MNIST, 
содержащий 70 тысяч полутоновых изображений размера $28 \times 28$ пикселей, 
представляющих рукописные цифры от 0 до 9~\cite{MNIST}. Данный набор разделен 
на две части:тренировочная выборка – 60 тысяч изображений,тестовая 
выборка – 10 тысяч изображений.

  Пример данных из базы MNIST представлен на рисунке~\ref{fig:mnist}.

\insertfigure{mnist}{pic/ex_mnist.png}{Пример данных из MNIST}{12cm}

Из тренировочного набора 
выделяется 1000 изображений для валидации, оставшиеся 59000 используются для 
непосредственного обучения модели. 

Перед подачей в нейросеть изображения проходят предварительную обработку, 
включающую нормализацию значений пикселей. Для упрощения 
процесса обучения выполняется нормализация данных таким образом, чтобы значение 
каждого пикселя было в диапазоне от $-1$ до $1$, а стандартное отклонение 
($\sigma$) составляло $0.5$. Это позволяет сделать градиентный спуск более 
устойчивым и ускорить процесс сходимости модели. Пример такой обработки 
представлен на рисунке~\ref{fig:norm-img}.

\insertfigure{norm-img}{pic/norm_img.png}{Пример обработанных данных из MNIST}{9cm}

Архитектура модели L2DST представляет собой последовательное соединение двух 
линейных слоёв с функцией активации гиперболический тангенс \(\tanh\) и одного 
полносвязного слоя с функцией активации softmax, применяемой для выполнения 
классификации. Основной особенностью реализации является использование слоёв 
LazyLinear, позволяющих автоматически определить размер входного пространства 
на этапе первого запуска модели.

Обучение модели осуществляется с использованием стохастического оптимизатора 
Adam, известного своей эффективностью на широком спектре задач глубокого 
обучения. Начальное значение скорости обучения задано равным 
\(1.5\times10^{-3}\).
Для обеспечения устойчивого процесса обучения применяется поэтапное снижение 
скорости обучения с помощью механизма StepLR. Каждые 10 эпох происходит 
уменьшение текущей скорости обучения на 3\%, что позволяет модели более точно 
подстраиваться к минимальному значению функции потерь на поздних этапах 
обучения.

На рисунке~\ref{fig:loss} показан график функции потерь за 370 эпохах.

\insertfigure{loss}{pic/loss.png}{Графики функции потерь на обучающем и тестовом наборе}{10cm}

После завершения каждой эпохи вычисляются средние значения ошибки на 
тренировочном и валидационном наборах данных.  
Python описание процесса обучения нейронной сети прямого распространения 
приведено в приложении Б.